%% Page 4: Figure 3 (one exemplar per task family), the models, and the first result (emergence).
\crumbs{§3.1 Setup \,›\, Figure~3: exemplars per task \,›\, Language models \,›\, §3.2 Results}
\insection{setup,results}

%% ---- marks on the original page
\badge{1}{35.5}{27.7}
\badge{2}{139.8}{27.7}
\badge{3}{35.5}{102.3}
\ring[ok]{115.6}{54.3}{121.1}{57.0}
\bracket{4}{154.3}{165.5}
\bracket{5}{167.7}{213.6}
\bracket{6}{219.9}{254.1}
\highlight[hlpink]{151.0}{239.0}{177.9}{242.6}
\highlight[hlpink]{38.1}{242.9}{90.5}{246.4}
\highlight[hlgreen]{163.9}{242.9}{177.8}{246.4}
\highlight[hlgreen]{38.1}{246.6}{64.5}{250.2}

\block[y=26]{%
\heading{Figure 3: one exemplar per task}
\begin{keyidea}
The format never changes, only what the reasoning talks about. The highlighted text in each box is the chain of thought.
\end{keyidea}
\num{1}\ \textbf{Arithmetic} (the only family used in §3): free response (the Figure 1 problem) and multiple choice:
$9+90(2)+401(3)=1392$ $\rightarrow$ (b) \green{(circled)}.}

\block{%
\num{2}\ \textbf{Commonsense}: CSQA (top right), StrategyQA, Date and Sports understanding (middle row). The chain
recalls a fact (a pear is about $0.6\ \mathrm{g/cm^3}$, less than water) and then concludes.}

\block[y=100]{%
\num{3}\ \textbf{Bottom row}:
\begin{points}
  \item SayCan (robot): an explanation, then a 5-step plan;
  \item symbolic: last letters of ``Lady Gaga'' $\rightarrow$ ``ya''; coin flip = keep track of the state
        (flipped once, odd $\rightarrow$ tails).
\end{points}
Commonsense is §4 and symbolic §5 of the paper; full prompts in Appendix G.}

\block[y=154]{%
\num{4}\ The \textbf{same 8 exemplars} for every maths benchmark except AQuA (multiple choice): 4 exemplars from its
training set (Appendix Table 21).}

\block[y=168]{%
\num{5}\ \textbf{Five model families}, several sizes each:
\begin{smalltable}{>{\bfseries}l l}
GPT-3 (InstructGPT) & 350M, 1.3B, 6.7B, 175B \\
LaMDA & 422M, 2B, 8B, 68B, 137B \\
PaLM & 8B, 62B, 540B \\
UL2 \quad Codex & 20B \quad code-davinci-002 \\
\end{smalltable}
\textbf{Greedy decoding} (always the most likely next token). LaMDA: average over 5 random orders of the
exemplars; one order for the others (little variance).}

\block[y=222]{%
\num{6}\ \textbf{Takeaway 1: an emergent ability.} Chain of thought \hl[hlpink]{does not help small models} and only
pays off with models of \hl[hlgreen]{about 100B parameters}.
\begin{warning}
Small models wrote fluent but illogical chains and did \textbf{worse} than standard prompting: this is not a trick
for small models.
\end{warning}}
